\documentclass[instructions.tex]{subfiles}
\begin{document}
\cleardoublepage
\markboth{NOTICE}{NOTICE}
\section*{NOTICE\\
\large SUR LA VIE DU P. HUMBERT.}
\addcontentsline{toc}{section}{NOTICE}

Hubert Humbert naquit en 1685 ou 1686, au village de Vanclans, paroisse de Nodz, département du Doubs. Ses parents étaient d’honnêtes cultivateurs ; chargés d’une nombreuse famille, ils ne négligèrent rien pour former leurs enfants à la vertu et à la piété chrétienne. Dieu les en récompensa ; ils virent trois de leurs fils honorés du sacerdoce. Celui qui est l’objet de cette notice était l’aîné. Il ne paraît pas qu’il ait reçu plus de soins pour sa première éducation que n’en reçoivent communément les enfants dans les petites écoles de campagne ; mais il se distingua dans ses cours d’humanités, de philosophie et de théologie. Il embrassait à la fois et avec succès plusieurs genres d’études : physique, histoire, poésie même ; et il eût pu n’être médiocre dans aucun. Sa piété et son zèle l’attachèrent de bonne heure à l’art oratoire, au ministère de la prédication. N’étant encore que sous-diacre, il fit dans une assemblée de prêtres et de jeunes ecclésiastiques, à Besançon, un discours où l’on vit dès-lors percer son génie et le talent rare qu’il montra ensuite pour la chaire.

Associé, en 1714, aux prêtres missionnaires de Besançon, dits de Beaupré, il se fit remarquer par la pénétration de son esprit, la solidité de son jugement et l’étendue de ses connaissances. Une mémoire heureuse, un travail facile, secondaien t ses autres dispositions. Il avait, dans ses discours et dans son débit, un genre particulier et qui n’était qu’à lui : il savait en faire sortir fréquemment des traits frappans d’éloquence ; la grandeur et la beauté des images s’y joignaient comme naturellement à la force des pensées et à la solidité des preuves. Sa voix, sans être forte, se faisait entendre dans les plus grands vaisseaux. Son air doux et majestueux commandait l’attention, et semblait persuader d’avance son auditoire.

Quoique par sa vocation il fût principalement destiné à prêcher la parole de Dieu dans les campagnes, il l’annonça cependant fréquemment dans les villes, en mission ou dans d’autres circonstances, et toujours à la grande satisfaction de ses auditeurs et de ceux en particulier qui étaient plus en état de le juger. Aussi produisait-il les plus grands fruits dans le saint ministère ; les fidèles de toutes les classes s’adressaient à lui avec la même confiance. Estimé, recherché, chéri de tous, il se communiquait très-peu, n’aimait à se prêter qu’aux besoins des âmes, et ne voulait avoir aucune assiduité chez les grands. S’il en voyait de temps en temps, c’était le zèle seul qui le conduisait dans leurs maisons, soit pour rétablir l’union ou la subordination dans les familles, soit pour ramener des jeunes gens dans les voies de la sagesse, soit enfin pour faire cesser quelques scandales.

Si l’on accourait de toutes parts à ses sermons, les personnes de tous les rangs n’étaient pas moins empressées de l’avoir pour guide de leur conscience au saint tribunal. Souvent il fut appelé aux retraites ecclésiastiques, et il y seconda avec succès les directeurs du séminaire. Il fut ainsi, pendant de longues années, l’apôtre des simples fidèles et des pasteurs ; et, comme il joignait toutes les vertus chrétiennes et sacerdotales au zèle et à la science, il eut le mérite si rare de servir long-temps de lumière et d’exemple à tout le diocèse. Il mourut en 1778, âgé de 92 ans, après avoir travaillé plus de 60 ans dans les missions, et passé 65 ans dans la société des missionnaires, dont il fut long-temps le directeur.

Le Père Humbert a laissé les ouvrages suivans :

1.º \textit{Pensées sur les plus importantes vérités de la Religion} ;

2.º \textit{Instructions chrétiennes pour les jeunes gens} : ces deux ouvrages sont justement estimés, et le premier est répandu partout ;

3.º \textit{Instructions abrégées sur les devoirs et les exercices du chrétien}, à l’usage des missions des prêtres missionnaires du diocèse de Besançon ;

4.º \textit{Instruction sur les égaremens de l’esprit et du cœur humain et sur les vertus nécessaires au salut} ;

5.º plusieurs \textit{Cantiques spirituels} insérés dans les recueils de Cantiques à l’usage du diocèse de Besançon.

Il a laissé aussi, en manuscrit, un ouvrage sur les vertus ecclésiastiques, les devoirs des pasteurs et le gouvernement des paroisses ; une histoire assez étendue de la communauté des missionnaires de Beaupré ; des sermons : mais il n’en reste qu’un bien petit nombre, les autres ayant été perdus. Dans tout ce qui reste du Père Humbert, on voit toujours l’homme éclairé, le prêtre vertueux et plein de zèle, et enfin le savant formé par sa longue expérience et par son application constante à l’étude.

\end{document}
